\chapter{Results}
\section{Physics-Based Weight Calculation}
    The foundational step of this project involved converting traditional analytical gearbox weight 
    calculations into a robust Python script. The specific architecture studied in this project is a 
    dual-stage helicopter main rotor gearbox. The powertrain consists of a high-speed input spiral 
    bevel gear stage (which changes the axis of rotation from horizontal to vertical) followed by a 
    low-speed, high-torque epicyclic planetary stage driving the main rotor mast.

    To accurately model the complex physics of a helicopter main rotor gearbox, the script accepts 31 
    independent input parameters. These inputs were chosen based on typical, realistic ranges utilized 
    in the modern aerospace industry. Key inputs include:
		\begin{itemize}
			\item Fundamental gear geometry (helix angle $\beta$, normal pressure angle $\alpha_n$, face-width to diameter ratio $\lambda_d$)
			\item Planetary configuration (number of planets $N_p$)
			\item Casing clearances
			\item Material properties
		\end{itemize}

    From these inputs, the script calculates the initial geometry of the entire gearbox assembly. 
    It sequentially evaluates the required physical dimensions for the input bevel gears, the planetary gear 
    stages, the shafts, and the outer housing required to safely and reliably transmit the power without yielding 
    contact fatigue or bending stresses. 

    A critical design parameter in dual-stage gearboxes is the intermediate shaft speed, which perfectly 
    balances the torque and weight distribution between the high-speed input bevel gears and the low-speed, 
    high-torque planetary stage. To find this optimum intermediate speed, a Sweep Search Algorithm was utilized.  

    To ensure the physical validity of the initial sizing calculations, the foundational gearbox configuration 
    was verified in Romax where the simplified analytical equations used in Python model accurately represent 
    sound aerospace designs.

\section{Sensitivity Study}
    A sensitivity study is used to quantify the absolute importance of design variables and understand their direct, 
    often non-linear, effect on system weight and physical size. The input variables were systematically varied to 
    study their effect within practical, manufacturable ranges. Through this sensitivity study, the primary driving 
    variables were identified that dictate most of the gearbox's mass as shown in \autoref{tab:Table2} and \autoref{tab:Table3}.

        \begin{table}[ht]
			\renewcommand{\arraystretch}{1.5}
			\centering
			\caption{Sensitivity Analysis Result for the weight}
			\label{tab:Table2}
			\small
			\begin{tabular}{|c|c|c|}
				\hline
        		\textbf{Rank} & \textbf{Variables} & \textbf{Delta} \\					
				\hline
				\textbf{1} & $\alpha_{n}$ & 1.8405 \\
				\hline
				\textbf{2} & P & 0.8069 \\
				\hline
				\textbf{3} & $\lambda_{d}$ & 0.6802 \\
				\hline
				\textbf{4} & $\beta$ & 0.5392 \\
				\hline
				\textbf{5} & $\rho_{shaft}$ & 0.5344 \\
				\hline
				\textbf{6} & $\sigma_{shaft}$ & -0.4344 \\
				\hline
				\textbf{7} & $n_{rotor}$ & -0.4324 \\
				\hline
				\textbf{8} & $N_{p}$ & -0.4109 \\
				\hline
				\textbf{9} & $K_{a}$ & 0.4064 \\
				\hline
				\textbf{10} & $K_{t}$ & 0.4064 \\
				\hline
            \end{tabular}
		\end{table}

        \begin{table}[ht]
			\renewcommand{\arraystretch}{1.5}
			\centering
			\caption{Sensitivity Analysis Result for the size}
			\label{tab:Table3}
			\small
			\begin{tabular}{|c|c|c|}
				\hline
        		\textbf{Rank} & \textbf{Variables} & \textbf{Delta} \\					
				\hline
				\textbf{1} & $\alpha_{n}$ & 2.4391 \\
				\hline
				\textbf{2} & $N_{p}$ & -1.9291 \\
				\hline
				\textbf{3} & $\beta$ & 0.7988 \\
				\hline
				\textbf{4} & P & 0.6978 \\
				\hline
				\textbf{5} & $\lambda_{d}$ & 0.6178 \\
				\hline
				\textbf{6} & $\sigma_{shaft}$ & -0.3770 \\
				\hline
				\textbf{7} & $C_{t}$ & 0.3015 \\
				\hline
				\textbf{8} & $L_{clearance}$ & 0.2250 \\
				\hline
				\textbf{9} & $\sigma_{\text{allowable contact}}$ & -0.2219 \\
				\hline
				\textbf{10} & $Z_{E}$ & 0.2210 \\
				\hline
            \end{tabular}
		\end{table}

\section{Dataset Generation and Material Selection}
    To train an Artificial Intelligence (AI) model capable of predicting complex physics, we require a 
    massive dataset. The top driving variables identified in the sensitivity study were varied to generate 
    tens of thousands of unique gearbox designs as shown in \autoref{fig:gearbox19}

    The samples for data generation were generated using Sobol Sequence Sampling rather than 
    traditional Monte Carlo random sampling. Monte Carlo methods tend to cluster data points randomly, 
    leaving massive \text{'gaps'} or \text{'holes'} in the multi-dimensional parameter space. Sobol sequences are 
    quasi-random, low-discrepancy sequences explicitly designed to distribute points as uniformly as 
    mathematically possible across an N-dimensional space. This ensures the ML model receives a perfectly 
    balanced learning of the edges of the design space just as well as the center. The dataset was generated 
    using 11 distinct material profiles such as: AISI 9310 Steel, Nitrided Steels, Ferrium C61 and Titanium.
        \begin{figure}[h]
            \centering
            \includegraphics[width=1\textwidth]{gearbox/gearbox19}
            \caption{Generated Dataset}
			\label{fig:gearbox19}
        \end{figure}
        \FloatBarrier 

\section{Machine Learning Models}
    With the data generated, multiple machine learning architectures were investigated to serve as 
    the surrogate model. The goal of a surrogate model is to predict the true physics calculations in 
    microseconds rather than seconds. Different models such as Random Forest, Multi-Layer Perceptron (MLP) 
    and HistGradientBoosting (HGBT) were evaluated. Model performance was measured using the Coefficient of 
    Determination ($R^2$). $R^2$ is a statistical measure that represents the proportion of the variance 
    for a dependent variable that's explained by the independent variables. An $R^2$ of 1.0 indicates 
    flawless, perfect prediction.

    Based on benchmarking across various maximum tree depths as shown in \autoref{tab:Table4}, 
    HistGradientBoosting consistently outperformed all other models, achieving a $R^2$ score of > 0.89 
    with a Mean Absolute Error (MAE) of less than 42 kg on a 500+ kg gearbox. Therefore, HistGradientBoosting 
    was selected as the core architecture.
        \begin{table}[ht]
			\renewcommand{\arraystretch}{1.5}
			\centering
			\caption{Different Model Performance Comparison}
			\label{tab:Table4}
			\small
			\begin{tabular}{|c|c|c|c|c|}
				\hline
        		\textbf{Model Name} & \textbf{Depth} & \textbf{$R^2$} & \textbf{Mean Error (kg)} & \textbf{Time} \\					
				\hline
				\textbf{Random Forest} & 10 & 0.8721 & 46.66 & 23.0s \\
				\hline
				\textbf{XGBoost} & 10 & 0.8886 & 42.46 & 45.4s \\
				\hline
				\textbf{HistGradientBoosting} & 10 & 41.85 & 10 & 10.9s \\
				\hline
				\textbf{Random Forest} & 25 & 0.8823 & 44.18 & 9.7s \\
				\hline
				\textbf{XGBoost} & 25 & 0.8626 & 46.80 & 324.5s \\
				\hline
				\textbf{HistGradientBoosting} & 25 & 41.88 & 10 & 4.9s \\
				\hline
            \end{tabular}
		\end{table}

\section{The Discontinuity Problem and Model Cascade Architecture}
    During initial training, a severe physical discontinuity problem was discovered. In real-world manufacturing, 
    gear teeth cannot be cut to any arbitrary decimal size; they must conform to standard tooling sizes known as 
    ISO Modules (e.g., 2.0 mm, 2.5 mm, 3.0 mm).  When the continuous physics math predicts a required gear module 
    of 2.1 mm, it must artificially snap up to the 2.5 mm standard. This snapping behavior creates massive, sudden 
    jumps in the final gearbox weight. Standard Machine Learning models fundamentally struggle to learn these 
    sudden discontinuous step-functions, leading to high error rates.

    Therefore, instead of training one model to do everything, the brain was splitted into two sequential models:
        \begin{enumerate}
            \item Model 1 (The Sizer): Predicts the continuous, exact mathematical module required by the physics formulas.
            \item The Snapper: Python function intercepts Model 1's prediction and perfectly snaps it to the nearest standard ISO module.
            \item Model 2 (The Weigher): Takes all the original geometric inputs plus the newly snapped module feature to predict the final Weight and Volume.
        \end{enumerate}

    This architecture bridges the gap between pure continuous mathematics and discrete manufacturing limitations, 
    eliminating the discontinuity error. This results in improvement in the accuracy of the AI model as shown in \autoref{tab:Table5}.
        \begin{table}[ht]
			\renewcommand{\arraystretch}{1.5}
			\centering
			\caption{Model Cascade Performance}
			\label{tab:Table5}
			\small
			\begin{tabular}{|l|c|c|}
				\hline
        		\textbf{Model Name} & \textbf{$R^2$} & \textbf{Mean Error} \\					
				\hline
				\textbf{Sizer} & 0.985 & 0.0003 mm \\
				\hline
				\textbf{Weigher} & 0.984 & 15.36 kg \\
				\hline
            \end{tabular}
		\end{table}

\section{Feature Importance (Permutation Importance)}
    Permutation Importance physically breaks the relationship between a specific feature and the target variable 
    by randomly shuffling that column's data. We then measure how severely the model's $R^2$ accuracy drops. 
    A massive drop in accuracy indicates that the model relies heavily on that feature.

    The results confirmed our physics models: the physical size of the gear teeth was the single most 
    critical feature, followed closely by the geometric knobs ($\beta$, `$\alpha_n$`, $\lambda_d$) and the planetary 
    configuration ($N_p$) as shown in \autoref{fig:gearbox20} and \autoref{fig:gearbox21}. These specific features therefore can be selected as the 
    primary variables for the final optimization phase.
        \begin{figure}[h]
            \centering
            \includegraphics[width=0.8\textwidth]{gearbox/gearbox20}
            \caption{Feature importance of the sizer}
			\label{fig:gearbox20}
        \end{figure}
        \FloatBarrier 

        \begin{figure}[h]
            \centering
            \includegraphics[width=0.8\textwidth]{gearbox/gearbox21}
            \caption{Feature importance of the weigher}
            \label{fig:gearbox21}
        \end{figure}
        \FloatBarrier 